(1) 理想气体 \(N\) 个粒子质量 \(m\)，限制在横截面 \(A\)、高度 \(H\) 的容器中，底部 \(z=0\)，重力势为 \(mgz\)，温度 \(T\) 均匀。构造粒子的空间配分因子及归一化高度密度。

(2) 求压强 \(p(z)\)，从积分精确证明底部与顶部压强差等于全部气体重量除以面积。

(3) 计算单粒子平均高度及总重力势能，并求低重力/矮柱极限和强重力/高柱极限。

(4) 求气体总平均能量，比较势能与动能，说明在等温环境下压强为何不按高度线性变化。
